\documentclass[10pt,a4paper]{article}

% ---------- Packages ----------
\usepackage[utf8]{inputenc}
\usepackage[T1]{fontenc}
\usepackage[english]{babel}

\usepackage[
    a4paper,
    margin=1.8cm,
    top=1.2cm,
    bottom=1.5cm,
    includehead,
    headheight=1.9cm,
    headsep=0.45cm
]{geometry}

\usepackage{graphicx}
\usepackage[table]{xcolor}
\usepackage{titlesec}
\usepackage{enumitem}
\usepackage{fancyhdr}
\usepackage{array}
\usepackage{tabularx}
\usepackage{xurl}
\usepackage{hyperref}

% ---------- Colors ----------
\definecolor{mainblue}{RGB}{25,55,110}
\definecolor{lightgray}{RGB}{245,245,245}

% ---------- Links ----------
\hypersetup{
    colorlinks=true,
    linkcolor=mainblue,
    urlcolor=mainblue,
    citecolor=mainblue
}

% ---------- Header ----------
\pagestyle{fancy}
\fancyhf{}
\renewcommand{\headrulewidth}{0.4pt}
\renewcommand{\footrulewidth}{0pt}

\fancyhead[L]{%
    \begin{minipage}[c]{0.32\textwidth}
        \raggedright
        \IfFileExists{uzi_logo.png}{\includegraphics[width=3.2cm]{\detokenize{uzi_logo.png}}}{\textbf{UZH}}
    \end{minipage}
}

\fancyhead[C]{%
    \begin{minipage}[c]{0.36\textwidth}
        \centering
        \textbf{\small MSc Thesis Proposal}\\[-0.1em]
        \scriptsize Department of Informatics
    \end{minipage}
}

\fancyhead[R]{}

\fancyfoot[C]{\thepage}

% ---------- Section Formatting ----------
\titleformat{\section}
  {\color{mainblue}\normalfont\large\bfseries}
  {}
  {0em}
  {}

\titlespacing*{\section}{0pt}{0.6em}{0.25em}

% ---------- Compact Lists ----------
\setlist[itemize]{noitemsep, topsep=2pt, leftmargin=1.2em}
\setlist[enumerate]{noitemsep, topsep=2pt, leftmargin=1.4em}

% ---------- Tables ----------
\renewcommand{\arraystretch}{1.08}
\setlength{\tabcolsep}{6pt}

\begin{document}

% ---------- Title ----------
\begin{center}
    {\Large \textbf{Standard versus Agentic Retrieval-Augmented Generation for\\[0.15em]Heterogeneous Industrial Documentation: A Query-Complexity Evaluation}}
\end{center}

\vspace{0.35em}

% ---------- Student Information ----------
\noindent
\begin{tabularx}{\textwidth}{>{\bfseries}l X >{\bfseries}l X}
Student Name: & Necati Furkan Colak       & Student ID:     & 24746323 \\
Program:      & MSc Computer Science      & Semester:       & Spring 2026 \\
Supervisor:   & [Supervisor Name]         & Co-Supervisor:  & [Co-Supervisor Name] \\
Email:        & necatifurkan.colak@uzh.ch & Date:           & \today \\
\end{tabularx}

\vspace{0.5em}
\hrule
\vspace{0.5em}

\section{Abstract}

Industrial knowledge is spread across manuals, procedures, standards, reports, and tables. Standard retrieval-augmented generation (RAG) \cite{lewis2020rag} answers simple questions with one retrieval step, but it fails when a question needs several documents combined, a query decomposed, or conflicting statements reconciled. Agentic RAG adds query planning, iterative retrieval, and answer verification \cite{yao2023react, singh2025agentic}, which raises token cost and latency on every call. This thesis compares a standard RAG pipeline and a bounded agentic RAG pipeline on the same corpus of heterogeneous industrial documents, with questions grouped by complexity: single-document, multi-document synthesis, and conflicting-evidence. It measures answer correctness, faithfulness, citation precision, and multi-hop completeness against latency, tokens, and cost per correct answer, with paired statistical tests per question group. The result identifies the query conditions under which agentic retrieval gives a measurable benefit, and those where it only adds cost, so that a system does not have to run the expensive path on every query.

\section{Background and Motivation}

An engineer asking a documentation system a question may want a single fact from one manual, or a value that only makes sense once three documents are read together, or a resolution between a general standard and a specific exception. Standard RAG treats all of these the same way: embed the query, retrieve once, generate. It handles the first case well and degrades on the others. Agentic RAG can plan the query, retrieve more than once, and check its own answer, but each of those steps costs additional model calls. Whether that cost buys enough quality depends on the question, and that dependence is what this thesis measures.

\section{Research Gap}

Agentic RAG is often described as more flexible than standard RAG \cite{singh2025agentic}, yet there is little controlled measurement of which query types actually benefit from the added planning, iteration, and verification, and which only pay extra cost and latency. Benchmarks show that standard pipelines lose accuracy on queries whose evidence spans several documents \cite{tang2024}. Interview studies with industrial adopters report that evaluation, not model capability, is the main obstacle to deployment \cite{ragindustry2025}. No published study conditions the standard-versus-agentic comparison on query complexity for heterogeneous industrial documentation.

\section{Research Objective}

Compare a standard RAG pipeline with a bounded agentic RAG pipeline on industrial document questions of different complexity levels, measuring quality and cost together per level.

\section{Research Questions}

\begin{enumerate}
    \item \textbf{RQ1 (primary):} Under which query conditions does agentic RAG provide measurable benefits over standard RAG for heterogeneous industrial documentation?
    \item \textbf{RQ2:} Does agentic retrieval help most on single-document, multi-document synthesis, or conflicting-evidence questions?
    \item \textbf{RQ3:} How much do iterative retrieval and answer verification improve faithfulness, and does the quality gain offset the latency and token-cost increase?
\end{enumerate}

\section{Proposed Methodology}

\begin{itemize}
    \item \textbf{Data:} A corpus of 150 to 300 open technical documents (manuals, safety instructions, standards excerpts, troubleshooting documents). A benchmark of about 120 questions is split into three groups: direct single-document questions, multi-document synthesis questions, and ambiguous or conflicting-evidence questions. Questions are drafted with an LLM and verified by hand, with gold answers and source passages. No proprietary data is required.
    \item \textbf{Standard RAG:} hybrid retrieval (BM25 and dense), a reranker, and a single generation step.
    \item \textbf{Agentic RAG:} query classification and planning, iterative retrieval, an evidence-sufficiency check, and answer verification. At most two or three logical agent roles are used; the design does not spawn many autonomous agents.
    \item \textbf{Baselines:} LLM only; standard RAG; agentic RAG. Ablations remove iterative retrieval and remove verification.
    \item \textbf{Evaluation:} Recall@k; context relevance; answer correctness; faithfulness \cite{es2024ragas}; citation precision; multi-hop completeness; hallucination rate; latency; token usage; cost per correct answer. Analysis uses paired comparisons within each question group and reports effect sizes.
\end{itemize}

\section{Expected Contribution}

\begin{enumerate}
    \item \textbf{Academic:} Experimental evidence on the query-complexity conditions under which agentic RAG is worth its cost, rather than a single system demonstration.
    \item \textbf{Technical:} A reproducible standard-versus-agentic RAG benchmark pipeline with a query taxonomy and ablations.
    \item \textbf{Practical:} Design guidance on when standard RAG is enough, so an organization does not run expensive agentic RAG on every query.
\end{enumerate}

\section{Scope and Limitations}

The thesis does not build a full enterprise platform, a web-search agent, autonomous workflow execution, or multi-department use. It covers document retrieval and grounded answer generation only. It uses one corpus, three question groups, and a bounded agentic design.

\section{Feasibility}

The data is public technical documentation, so no data owner is needed. The agentic architecture is limited to two added functions beyond standard RAG (iterative retrieval and verification), which keeps implementation and ablation within a single-student budget. The experimental load is about 120 questions across three to five system variants.

\section{Preliminary Work Plan}

\begin{tabularx}{\textwidth}{|p{1.9cm}|X|}
\hline
\rowcolor{lightgray}
\textbf{Period} & \textbf{Planned Work} \\
\hline
Month 1 & Corpus assembly and query taxonomy. \\
\hline
Month 2 & Ground-truth benchmark construction with human verification. \\
\hline
Month 3 & Standard RAG implementation. \\
\hline
Month 4 & Agentic RAG implementation. \\
\hline
Month 5 & Quantitative experiments and ablations. \\
\hline
Month 6 & Analysis; thesis writing; final presentation. \\
\hline
\end{tabularx}

\begingroup
\renewcommand{\refname}{\texorpdfstring{\color{mainblue}}{}Preliminary References}
\scriptsize
\begin{thebibliography}{9}
\setlength{\itemsep}{0pt}

\bibitem{lewis2020rag}
P.\ Lewis et al.
\textit{Retrieval-Augmented Generation for Knowledge-Intensive NLP Tasks}.
NeurIPS, 2020.

\bibitem{yao2023react}
S.\ Yao et al.
\textit{ReAct: Synergizing Reasoning and Acting in Language Models}.
ICLR, 2023.

\bibitem{singh2025agentic}
A.\ Singh et al.
\textit{Agentic Retrieval-Augmented Generation: A Survey on Agentic RAG}.
arXiv:2501.09136, 2025.

\bibitem{tang2024}
Y.\ Tang and Y.\ Yang.
\textit{MultiHop-RAG: Benchmarking Retrieval-Augmented Generation for Multi-Hop Queries}.
arXiv:2401.15391, 2024.

\bibitem{es2024ragas}
S.\ Es et al.
\textit{RAGAS: Automated Evaluation of Retrieval Augmented Generation}.
EACL, 2024.

\bibitem{ragindustry2025}
\textit{Retrieval-Augmented Generation in Industry: An Interview Study on Use Cases, Requirements, Challenges, and Evaluation}.
arXiv:2508.14066, 2025.

\end{thebibliography}
\endgroup

\vspace{0.1em}

\noindent
\begin{tabularx}{\textwidth}{X X}
\textbf{Student Signature:} & \textbf{Supervisor Signature:} \\[0.9em]
\rule{0.42\textwidth}{0.4pt} & \rule{0.42\textwidth}{0.4pt} \\
\end{tabularx}

\end{document}
